\documentstyle[%


righttag%


,11pt%


,amssymb%


,russian%               remove in Eng.


,izvru%                 Set izven% in Eng.


% ,izvru-ks             for brief communications


]{amsart}





%       Math definitions





\newcommand{\tts}[1]{}





%\hoffset=-15mm%                  For Eng.


\setcounter{page}{3}


\god{2000}


\nomer{3~(454)} %                        Remove in Eng.


%\nomer{Vol.\,44, No.\,3}%               Set in Eng.


%\str{\pageref{begin}--\pageref{end}}%  Set in Eng.


\udk{519.214}





\author{\it О.В.\,Антонова, С.А.\,Бронникова, В.В.\,Давыдова, М.\,Ордоньез Кабрера}


% NB:   ^^ remove in Eng.


\title{О слабом законе больших чисел в банаховых


пространствах мартингального типа $\size{14}{18pt}\selectfont p$\\


при общем условии типа Чезаро}





\date{}


%\pagestyle{myheadings}%         Set in Eng.


\pagestyle{plain} %              Remove in Eng.


\begin{document}


%\thispagestyle{plain}%          Set in Eng.


\maketitle


\label{begin}





\section*{Введение}





Cтатья посвящена доказательству слабого закона больших чисел (СЗБЧ)


общего вида


$$


S_n= \sum^{N_n}_{j=1} a_{nj}


(V_{nj}-c_{nj}) @>{P}>> 0 \quad \text{при} \quad


n\to\infty,


$$


где


$\{a_{nj},\ j \geq 1,\ n \geq1\}$ --- схема серий вещественных чисел, 


$\{c_{nj},\ j \geq 1,\ n \geq 1\}$ --- так называемая ``центрирующая'' схема 


серий, составленная из специально подобранных условных ожиданий,


$\{N_n,\ n\geq 1\}$ --- последовательность целочисленных случайных величин и 


$\{V_{nj},\ j \geq 1,\ n \geq 1\}$ --- схема серий случайных элементов,


определенных на вероятностном пространстве $(\Omega, {\cal F},P)$


и принимающих значения в сепарабельном банаховом пространстве


${\cal X}$ с нормой $\|\cdot\|$. Случайный элемент $S_n$ называется


{\it взвешенной суммой} со {\it взвешенными коэффициентами}


$$


\{a_{nj},\ j \geq 1,\ n \geq 1\}.


$$





Предполагается, что банахово пространство


${\cal X}$ имеет мартингальный тип $p$, т.\,е. существует такая постоянная


$C$, что для всех мартингалов $\{S_n,\ n \geq 1\}$ со значениями в ${\cal X}$


выполняется неравенство


$$ 


\sup_{n \geq 1} E\|S_n\|^p \leq C \sum^\infty_{n=1} E\|S_n-S_{n-1}\|^p,


$$


где $S_0 \equiv 0$.





Заметим, что любое вещественное сепарабельное банахово пространство имеет


мартингальный тип $1$, тогда как пространства $L_p$ и $l_p$ ($1\leq


p<\infty$) имеют мартингальный тип $\min \{p,2\}$. Хорошо известно, что


если банахово пространство имеет мартингальный тип $p$, то оно имеет и


радемахеровский тип $p$. Однако понятие мартингального типа $p$ только


поверхностно схоже с понятием радемахеровского типа $p$. Более подробное


обсуждение можно найти в [1].





Основной результат данной работы --- теорема 3. Она является расширением


и обобщением на случай банахового пространства мартингального


типа $p$ и взвешенных сумм утверждений, доказанных в [2]--[7],


которые были установлены в схеме серий вещественнозначных


случайных величин. Заметим, что СЗБЧ доказывается в предположении


выполнимости условия типа Чезаро, которое более общо, чем условие из [4].





Вопрос о выполнимости СЗБЧ при условии типа Чезаро был впервые поставлен в [2].


Далее тот же вопрос рассматривался в [4], но с изменением условия типа Чезаро 


на условие типа Чезаро--Хонга. Причем оба результата были установлены в схеме


серий вещественнозначных случайных величин.


Существенное продвижение в этом направлении было осуществлено в [8], где


результаты


[2] и [4] были перенесены на случай банаховых пространств мартингального типа 


и неслучайных взвешенных сумм. Следующие работы улучшали полученные результаты.


Так, например, в [3] СЗБЧ рассматривался для случайных взвешенных сумм, в [9]


меняются условия на взвешенные коэффициенты, [10] обобщает [3] на случай


банахового пространства мартингального типа.





В данной статье найдено универсальное условие типа Чезаро. Оно более общее,


чем условие Чезаро--Хонга и позволяет получить все предыдущие результаты


как простые следствия. При этом ослабляется условие (4) из теоремы статьи [9].


Кроме этого, накладываются условия на скорость возрастания


взвешенных коэффициентов $\{a_{nj},\ j \geq 1,\ n \geq 1\}$ и на 


маргинальные распределения случайных величин


$\{\|V_{nj}\|,\ j \geq 1,\ n \geq 1\}$,


тем самым получается другой, более подходящий для решения конкретных задач,


закон больших чисел. Следует отметить, что все исследования по этой тематике


вдохновлены замечательной статьей Аллана Гута [2].





\section*{Основной результат}





Прежде чем сформулировать основной результат, приведем две


теоремы, которые доказаны в [10] и [9] соответственно.





\begin{thm}[{\series{m}\shape{n}\selectfont\cite{7}}]


$\mathrm I$. Пусть $\{V_{nj},\ j\geq 1,\ n \geq 1\}$ ---


схема серий случайных элементов в вещественном сепарабельном банаховом


пространстве мартингального типа $p$ $(1 \leq p


\leq 2)$, $\{N_n,\ n\geq 1\}$ --- последовательность целочисленных


случайных величин и $k_n\ (\rightarrow \infty)$ --- некоторая


неслучайная последовательность натуральных чисел, для которых


\begin{equation}


P\{N_n> k_n\}= o(1)\quad \text{при}\quad n\to\infty .


\end{equation}


Пусть схема серий вещественных чисел $\{a_{nj},\ j \geq 1,\ n \geq 1\}$


и последовательность


$\{f(n),\ n\geq 1\}$, определенная как $f(n)= 1/ \max\limits_{1\leq


j\leq k_n} |a_{nj}|$, удовлетворяют условию


\begin{equation}


k_n f^{-p}(n)= o(1)\quad\text{при}\quad n\to\infty .


\end{equation}





$\mathrm{II}$. Пусть $\{g(m),\ m\geq 0\}$ --- такая неубывающая 


последовательность


неотрицательных чисел, что выполняется


\begin{equation}


\sum_{m=1}^{k_n-1} \frac{g^{p}(m+1) -


g^{p}(m)}{m}= O(f^{p}(n)/k_n)\quad \text{при}\quad n\to\infty ,


\end{equation}


а также равномерное условие типа Чезаро


\begin{equation}


\lim_{m \rightarrow \infty}\, \sup_{n \geq 1}\,


\frac{1}{k_n}\sum^{k_n}_{j=1} m P\{\|V_{nj}\|> g(m)\} = 0.


\end{equation}





Тогда имеет место СЗБЧ


\begin{equation}


\sum^{N_{n}}_{j=1} a_{nj}(V_{nj}-E(V'_{nj}\mid{\cal F}_{n,j-1}))


@>{P}>> 0\quad \text{при}\quad n\to\infty,


\end{equation}


где $V'_{nj} = V_{nj}I(\|V_{nj}\|\leq g(k_n))$,


${\cal F}_{nj} = \sigma(V_{ni},\ 1 \leq i \leq j)$, $j \geq 1$, $n \geq 1$,


и ${\cal F}_{n0} = \{\emptyset, \Omega\}$, $n \geq 1$.


\end{thm}





\begin{thm}[{\series{m}\shape{n}\selectfont\cite{6}}]


Пусть выполняются все условия теоремы $1$


за исключением $(3)$, которое заменено на


\begin{equation}


\sum_{m=1}^{g(k_n)-1} \frac{m^{p-1}}{h(m)}


= O(f^p(n)/k_n)\quad \text{при}\quad n\to\infty,


\end{equation}


где


$h(m)=\min\{n\ge 0: g(n)\ge m\}$ --- обратная последовательность к $g(m)$.


Тогда имеет место $(5)$.


\end{thm}





Теперь сформулируем результат, из которого обе


теоремы могут быть получены как простые следствия. 





\begin{thm}


Пусть выполняется часть $\mathrm I$ теоремы $1$.


Пусть $\{\beta(m),\ m\geq 1\}$, $\{\alpha(m),\ m\geq 1\}$ и


$\{\wt g(m),\ m\geq0\}$


--- неубывающие последовательности положительных чисел и


$\{m_n,\ n\geq 1\}$ --- последовательность натуральных чисел,


для которых


\begin{equation}


\beta(m_n)= \wt g(n), \quad


\alpha(m_n)\ge k_n


\end{equation}


и


\begin{equation}


\sum_{m=1}^{m_n-1} \frac{\beta^{p}(m+1) - \beta^{p}(m)}{\alpha(m)}=


O(f^{p}(n)/k_n)\quad \text{при}\quad n\to\infty.


\end{equation}


Предположим, что выполняется равномерное условие типа Чезаро


\begin{equation}


\lim_{m \rightarrow \infty}\, \sup_{n \geq 1}\, \frac{1}{k_n}


\sum^{k_n}_{j=1} \alpha(m) P\{\|V_{nj}\|> \beta(m)\} = 0.


\end{equation}


Тогда имеет место СЗБЧ $(5)$, где


$V'_{nj} = V_{nj}I(\|V_{nj}\|\leq \wt g(n))$.


\end{thm}





\begin{pf}


Из условий (1), (6), (7)


и (9) следует, что для любого $\varepsilon > 0$


и $n \geq 1$


\begin{multline*}


P\bigg\{\bigg\|\sum^{N_{n}}_{j=1} a_{nj}V_{nj} - 


\sum^{N_{n}}_{j=1}


a_{nj}V'_{nj}\bigg\| > \varepsilon \bigg\}


\leq P\{N_n> k_n\} +


P\bigg\{\bigcup^{k_{n}}_{j=1} [V_{nj} \neq 


V'_{nj}]\bigg\} \le \\


%


\leq o(1) +


\frac{1}{k_{n}} \sum^{k_{n}}_{j=1} k_n P\{\|V_{nj}\|


> \wt g(n)\}\leq o(1) + \frac{1}{k_n}


\sum^{k_n}_{j=1} \alpha(m_n) P\{\|V_{nj}\|> \beta(m_n)\}


= o(1).


\end{multline*}


Поэтому


$\sum\limits^{N_{n}}_{j=1} a_{nj} V_{nj} - \sum\limits^{N_{n}}_{j=1}


a_{nj}V'_{nj} @>{P}>> 0$


и достаточно доказать, что


$\sum\limits^{N_{n}}_{j=1} a_{nj}V'_{nj} - \sum\limits^{N_{n}}_{j=1}


a_{nj}c_{nj} @>{P}>>0$,


где $c_{nj}=E(V'_{nj}\mid{\cal F}_{n,j-1})$.





Для $n\ge 1$ и $m\ge 1$ обозначим


$B_m^n=\Big\{\|\sum\limits^{m}_{j=1} a_{nj}V'_{nj} - \sum\limits^{m}_{j=1}


a_{nj}c_{nj}\|>\varepsilon \Big\}$


и $D_n= \bigcup^{k_n}_{m=1} B_m^n$. Тогда по (1)


$$


P\{B_{N_n}^n\}\le P\{B_{N_n}^n,\ N_n\le k_n\} + P\{N_n> k_n\}\le


P\{D_n\} + o(1)


$$


и осталось доказать, что $P\{D_n\}= {\it o}(1)$.





Отметим, что по неравенству Маркова


$$


P\{D_n\}= P\bigg\{\max_{1\le m\le k_n}\bigg\|\sum^{m}_{j=1}a_{nj}


(V'_{nj} - c_{nj})\bigg\|> \varepsilon \bigg\}


\le C E\max_{1\le m\le k_n}\bigg\|\sum^{m}_{j=1}


a_{nj}(V'_{nj}-c_{nj})\bigg\|^p.


$$





Так как рассматриваемое банахово пространство имеет мартингальный тип $p$, то


$$


P\{D_n\}


\leq C \sum^{k_{n}}_{j=1}E(|a_{nj}|\, \|V'_{nj}-c_{nj}\|)^p


\leq C2^{p-1}f^{-p}(n)\sum^{k_{n}}_{j=1}E(\|V'_{nj}\|^p +


E\|c_{nj}\|^p)


\leq C 2^p f^{-p}(n) \sum^{k_{n}}_{j=1}E\|V'_{nj}\|^p


$$


(во второй оценке применяем неравенство Йенсена для условных


ожиданий (напр., [11], c.\,209) для выпуклой функции $|\cdot|^p)$.


Более того,


\begin{gather*}


\sum^{k_{n}}_{j=1}E\|V'_{nj}\|^p=


\sum^{k_{n}}_{j=1}\sum^{m_n}_{m=1}E\|V_{nj}\|^p I(\beta(m-1)<


\|V_{nj}\|\leq \beta(m))\le \\


%


\leq \sum^{k_{n}}_{j=1}\sum^{m_n}_{m=1}


\beta^p(m)(P\{\|V_{nj}\| >m-1\}


- P\{\|V_{nj}\|> m\})= \\


%


= \sum^{k_{n}}_{j=1}[\beta^p(1) P\{\|V_{nj}\| > 0\} -


\beta^p(m_n) P\{\|V_{nj}\|> \beta(m_n)\} + \\


%


+\sum^{m_{n}-1}_{m=1}(\beta^p(m+1)- \beta^p(m))


P\{\|V_{nj}\|>\beta(m)\}]\le \\


%


\leq \beta^p(1)k_n +k_n


\sum^{m_n-1}_{m=1}\frac{\beta^p(m+1)-\beta^p(m)}{\alpha(m)}b_{nm},


\end{gather*}


где


$b_{nm}=\frac{1}{k_n}\sum\limits_{j=1}^{k_n}


\alpha(m) P\{\|V_{nj}\|> \beta(m)\}$.


Тогда по равномерному условию типа Чезаро (9) получим


$\sup\limits_{n\ge 1}


b_{nm}= o(1)$ при $m\to\infty$.


Далее, по (2), (8) и лемме Теплица (см., напр., [12],


с.\,250) имеем $P\{D_n\}=o(1)$,


чем завершается доказательство теоремы 3.


\end{pf}





\begin{cor}


Если $\wt g(n)= g(k_n)$, $\beta(m)= g(m)$, $\alpha(m)=


m$ в теореме 3, то справедливо утверждение теоремы 1.


\end{cor}





При $\wt g(n)= g(k_n)$, $\beta(m)= m$ и


$\alpha(m)= h(m)$ в теореме 3 получим





\begin{cor}


Пусть в условиях теоремы 1 условие (3) заменено на (6),


где $h(m)=\min\{n\ge 0: g(n)\ge m\}$ --- обратная к $g(m)$


последовательность,


а условие (4) заменено на


\begin{equation}


\lim_{m \rightarrow \infty}\, \sup_{n \geq 1}\, \frac{1}{k_n}


\sum^{k_n}_{j=1} h(m) P\{\|V_{nj}\|> m\} = 0.


\end{equation}


Тогда имеет место (5).


\end{cor}








Заметим, что можно ослабить условие (4) теоремы 2, заменив его на условие 


(10). Доказательство того, что (10) слабее условия (4),


приводится в [9].





\begin{thebibliography}{99}





\bibitem{10}


Pisier G. {\em Probabilistic methods in the geometry of Banach


spaces} // Lect. Notes Math. -- 1986. -- V.\,1206. -- P.\,167--241.


\bibitem{3}


Gut A. {\em The weak law of large numbers for arrays} // Statist.


Probab. Lett. -- 1992. -- V.\,14. -- P.\,49--52.


\bibitem{4}


Hong D.H. {\em On the weak law of large numbers for


randomlyindexed partial sums for arrays} // Statist. Probab. Lett. 


-- 1996. -- V.\,28. -- P.\,127--130.


\bibitem{5}


Hong D.H., Oh K.S. {\em On the weak law of large


numbers for arrays} // Statist. Probab. Lett. -- 1995. -- V.\,22. --


P.\,55--57.


\bibitem{8}


Hong D.H., Sung S.H., Volodin A.I. {\em On the weak law


for randomly indexed partial sums for arrays} // Preprint of Pai


Chai University. Submitted to Statist. Probab. Lett. -- 1998. -- 7 p.


\bibitem{8}


Kovalski P., Rychlik Z. {\em On the weak law of


large numbers for randomly idexed partial sums for arrays} // Ann.


Univ. Mariae Curie-Sklodovska. Sect.~A. -- 1997. -- V.\,LI. -- \No\,1.


-- P.\,109--119.


\bibitem{11}


Sung S.H. {\em Weak law of large numbers for arrays} //


Statist. Probab. Lett. -- 1998. -- V.\,38. -- P.\,101--105.


\bibitem{1}


Adler A., Rosalsky A., Volodin A.I. {\em A


meanconvergence theorem and weak law for arrays of random elements


inmartingale type $p$ Banach spaces} // Statist. Probab. Lett. -- 1977. --


V.\,32. -- P.\,167--174.


\bibitem{6}


Hong D.H., Ord\'o\~nez Cabrera M., Sung S.H., Volodin A.I.


{\em Again on the weak law for randomly indexed partial sums for


arrays of random elements in martingale type $p$ Banach spaces} //


Extracto Math. -- 1999. -- V.\,14. -- \No\,1. -- P.\,45--50.


\bibitem{7}


Hong D.H., Ord\'o\~nez Cabrera M., Sung S.H., Volodin A.I.


{\em On the weak law for randomly indexed partial sums for


arrays of random elements in martingale type $p$ Banach spaces} //


Statist. Probab. Lett. -- 1999. -- P.\,177-185.


\bibitem{2}


Chow Y.S. Teicher H. {\em Probability


theory$:$ independence, interchangeability, martingales}. --


N.~Y.: Springer-Verlag, 1997.


\bibitem{8}


Лоэв M. {\em Теория вероятностей}. -- M.:


Ин. лит., 1962. -- 720 c.





\end{thebibliography}





\vskip 2cm





\post{\it Казанский государственный университет}{\it Поступила}


\post{\it Университет Севильи $($Испания$)$}{23.04.1999}





\label{end}


\end{document}








